\chapter{Conclusion}
\label{ch:conclusion}

This chapter summarizes the contributions of the RCC agent-based model with glucose sensing, discusses its current limitations, and outlines directions for future work.

% ======================================================================
\section{Summary of Contributions}
\label{sec:contributions}

This project extends the foundation of computational cancer immunology through systematic integration of metabolic dynamics into agent-based tumor microenvironment simulation.

\textbf{Chapter 1} established clinical and computational motivation for agent-based modeling in RCC research, demonstrating that computational modeling can address clinically relevant questions about treatment optimization impractical to explore through clinical trials alone.

\textbf{Chapter 2} systematically compared agent-based modeling with traditional differential equation approaches, provided technical justification for the Mesa to Repast4Py framework migration, and established the biological foundation for modeling 18 distinct cell types and their interactions.

\textbf{Chapter 3} presented the complete agent-based architecture demonstrating how complex biological interactions can be systematically decomposed into local agent behaviors that produce emergent system-level phenomena.

\textbf{Chapter 4} introduced the foundational metabolic layer with a continuous 3D glucose field overlaid on discrete agent grid, implementing the Warburg effect with tumor cells consuming glucose at 200× normal rates and enabling bidirectional coupling between glucose availability and immune function.

\textbf{Chapter 5} detailed the 16-gene representation enabling modeling of tumor clonal evolution and immune evasion, with codon-level gene representation providing unprecedented molecular detail in agent-based tumor modeling.

\textbf{Chapter 6} implemented the pharmacological intervention layer modeling ICI and TKI therapies individually and in combination, with modular drug architecture allowing systematic exploration of dosing schedules and combination ratios.

\textbf{Chapter 7} addressed sex disparity in RCC through explicit modeling of estrogen, testosterone, and progesterone effects on immune cell behavior, with biphasic estrogen response curve capturing hormonal complexity underlying sex-dependent cancer outcomes.

\textbf{Chapter 8} documented software engineering challenges of porting from Mesa to Repast4Py while preserving biological fidelity, developing custom spatial indexing system and grid utility functions.

\textbf{Chapter 9} presented the 92-parameter system balancing biological interpretability with optimization tractability, with hierarchical parameter architecture separating patient-specific clinical variables from system-level behavioral weights.

\textbf{Chapter 10} validated the integrated model against ARON-1 clinical data, demonstrating that the glucose-extended model can reproduce observed treatment outcomes with sex-stratified results showing females achieving 65% survival with ICI+TKI combination versus males at 55%.

% ======================================================================
\section{Model Limitations}
\label{sec:limitations}

\begin{itemize}
    \item \textbf{Single-rank execution} — Model operates on single MPI rank, glucose field requires halo exchange for distribution across ranks
    \item \textbf{Simplified glucose dynamics} — Models only glucose with linear diffusion, missing oxygen, lactate, and amino acid coupling
    \item \textbf{Fixed blood vessel topology} — Blood vessels inject glucose at fixed rate regardless of surrounding conditions
    \item \textbf{No spatial immune infiltration patterns} — All immune cells initialized randomly without modeling vascular entry patterns  
    \item \textbf{Simplified TCR matching} — Uses discrete threshold-based comparison rather than continuous binding affinities
    \item \textbf{Deferred multi-rank validation} — SharedContext infrastructure in place but not validated in distributed configurations
\end{itemize}

% ======================================================================
\section{Future Work}
\label{sec:future-work}

\textbf{Near-term (6-12 months):}
\begin{itemize}
    \item Multi-rank MPI scaling with glucose field distribution using domain decomposition and halo exchange
    \item Multi-metabolite system adding oxygen, lactate, and amino acids with coupled reaction terms
    \item Dynamic angiogenesis with blood vessel pruning and sprouting driven by VEGF gradients
\end{itemize>

\textbf{Medium-term (1-2 years):}
\begin{itemize}
    \item Clonal evolution tracking with lineage relationships and spatial competition modeling
    \item Patient-specific simulation integration using CT/MRI imaging data for personalized modeling
    \item GPU acceleration using CuPy or JAX for real-time simulation capabilities
\end{itemize>

\textbf{Long-term (2-5 years):}
\begin{itemize}
    \item Prospective clinical trial integration with regulatory approval for AI-assisted treatment planning
    \item Multi-cancer framework adaptation to melanoma, lung cancer, and hepatocellular carcinoma
    \item Combination therapy optimization platform for screening thousands of potential combinations in silico
    \item Biomarker discovery through computational screening of high-dimensional simulation output data
\end{itemize>

% ======================================================================
\section{Clinical Impact}
\label{sec:clinical-impact}

This work represents a fundamental shift toward quantitative, predictive approaches to cancer treatment that could transform clinical practice. The integration of metabolic dynamics with immune system modeling addresses why seemingly similar patients respond differently to identical treatments, moving beyond population averages toward genuinely personalized cancer treatment. The renal cell carcinoma agent-based model with glucose sensing demonstrates that mechanistic biological modeling can generate clinically relevant predictions about treatment outcomes, providing a foundation for computational precision medicine.